\documentclass[UTF8,a4paper,11pt,fontset=fandol]{ctexart}

\usepackage{geometry}
\usepackage{booktabs}
\usepackage{longtable}
\usepackage{array}
\usepackage{tabularx}
\usepackage{xcolor}
\usepackage{graphicx}
\usepackage{tikz}
\usetikzlibrary{arrows.meta,backgrounds,calc,fit,positioning,shapes.geometric}
\usepackage{enumitem}
\usepackage{fancyhdr}
\usepackage{listings}
\usepackage{hyperref}
\usepackage{caption}
\usepackage{float}
\usepackage{titlesec}
\usepackage{amsmath}

\geometry{left=2.5cm,right=2.5cm,top=2.6cm,bottom=2.6cm}
\setlength{\parindent}{2em}
\setlength{\parskip}{0.25em}
\linespread{1.18}

\definecolor{ztblue}{RGB}{36,87,145}
\definecolor{ztgreen}{RGB}{33,130,96}
\definecolor{ztorange}{RGB}{206,116,31}
\definecolor{ztred}{RGB}{176,61,54}
\definecolor{ztgray}{RGB}{245,247,250}
\definecolor{ztline}{RGB}{185,190,198}
\definecolor{codebg}{RGB}{248,248,248}

\titleformat{\section}{\Large\bfseries}{\thesection}{0.7em}{}
\titleformat{\subsection}{\large\bfseries}{\thesubsection}{0.6em}{}
\titleformat{\subsubsection}{\normalsize\bfseries}{\thesubsubsection}{0.5em}{}

\hypersetup{
  hidelinks,
  pdftitle={zeroTrace 项目文档},
  pdfauthor={动感光波biubiubiu}
}

\pagestyle{fancy}
\fancyhf{}
\lhead{zeroTrace 项目文档}
\rhead{OS 功能挑战赛 · proj40}
\cfoot{\thepage}
\setlength{\headheight}{14pt}
\renewcommand{\headrulewidth}{0.4pt}
\renewcommand{\footrulewidth}{0pt}

\lstdefinestyle{ztcode}{
  basicstyle=\ttfamily\small,
  backgroundcolor=\color{codebg},
  frame=single,
  rulecolor=\color{ztline},
  breaklines=true,
  columns=fullflexible,
  keepspaces=true,
  showstringspaces=false,
  xleftmargin=0.5em,
  xrightmargin=0.5em,
  aboveskip=0.8em,
  belowskip=0.8em
}
\lstset{style=ztcode}

\newcolumntype{Y}{>{\raggedright\arraybackslash}X}
\newcolumntype{P}[1]{>{\raggedright\arraybackslash}p{#1}}
\renewcommand{\arraystretch}{1.22}
\newcommand{\code}[1]{\texttt{#1}}
\newcommand{\doneexp}{\textcolor{ztgreen}{已通过}}

\begin{document}

\hypersetup{pageanchor=false}

\begin{titlepage}
  \centering
  \vspace*{2.0cm}
  {\zihao{2}\bfseries zeroTrace：轻量级用户态动态探针项目文档\par}
  \vspace{0.8cm}
  {\zihao{4}2026 年全国大学生计算机系统能力大赛操作系统设计赛\par}
  \vspace{0.2cm}
  {\zihao{4}OS 功能挑战赛 · proj40\par}
  \vspace{1.8cm}

  \renewcommand{\arraystretch}{1.35}
  \begin{tabular}{P{3.5cm}P{9.2cm}}
    \toprule
    参赛赛道 & 2026 年全国大学生计算机系统能力大赛操作系统设计赛 -- OS 功能挑战赛 \\
    赛题编号 & proj40 \\
    赛题名称 & 轻量级用户态动态探针设计与实现 \\
    项目名称 & zeroTrace \\
    队伍名称 & 动感光波biubiubiu \\
    学校 & 中国矿业大学 \\
    队员 & 舒钰尧、曹佳一、倪晨曦 \\
    指导教师 & 林果园 \\
    \bottomrule
  \end{tabular}

  \vfill
  {\large 项目文档\par}
  \vspace{0.2cm}
  {\large \today\par}
\end{titlepage}

\hypersetup{pageanchor=true}
\tableofcontents
\newpage

\begin{abstract}
zeroTrace 是一套面向 Linux 用户态程序的轻量级动态探针系统。传统 uprobe 在命中用户态函数时通常要切到内核侧处理事件，高频函数上的上下文切换会明显放大观测成本。zeroTrace 把安装阶段和命中阶段分开：tracer 进程负责安装、更新和清理 probe；函数命中后的快速路径留在目标进程用户态内部，由入口 patch、trampoline、payload stub 和共享 trace buffer 完成参数采集、返回值捕获和事件输出。

系统支持动态库函数追踪、前 6 个参数捕获、return probe、多 probe 共存、动态开关、线程组级 patch 安全、条件过滤、探针热更新和探针内调用目标进程函数，并提供 x86\_64 与 aarch64 两套 ISA 后端。日志沿用 perf/ftrace 风格字段，记录 comm、pid/tid、CPU、\code{CLOCK\_MONOTONIC} 时间戳和事件类型，便于和系统级 trace 事件放到同一条时间线上分析。本次 x86\_64 评测中，zeroTrace 单次 probe 额外开销均值为 174.45 ns，比 kernel uprobe 低 11.08 倍；在 Google Cloud T2A aarch64 环境中，zeroTrace 额外开销均值为 237.02 ns，比 kernel uprobe 低 2.02 倍。两套环境中的 probe 安装/卸载延迟均低于赛题指标要求。

\noindent\textbf{关键词：}用户态动态探针；ptrace；trampoline；return probe；trace buffer；多线程安全；热更新
\end{abstract}

\section{引言}

动态追踪的价值在于“程序仍然按原来的方式运行，但开发者能看到更细的执行事实”。在调试线上服务、分析性能抖动、定位库函数调用链时，开发者往往希望临时插入 probe，观察某个函数的参数、返回值和调用频率。Linux 已经提供 uprobe、perf、ftrace 等工具，但其中用户态函数 probe 多数需要在命中时进入内核。低频函数还能接受这种成本；一旦函数被百万次调用，观测本身就可能改变被观测对象。

proj40 要求实现一种轻量级用户态动态探针。普通调试器式做法可以很快 trace 到函数，但每次命中都触发内核处理，热点函数上的成本会被放大。zeroTrace 把这部分成本放到安装阶段：先使用 ptrace 取得目标进程控制权，注入 payload 并改写函数入口；目标进程恢复运行后，函数直接跳入进程内的 trampoline 和 stub，事件写入共享 ring buffer，再由 tracer 异步读取。

本文档整理 zeroTrace 的设计、实现和实验记录。后文会对照赛题要求说明支持情况，并依次介绍总体架构、关键实现、正确性测试和性能评测。

\subsection{设计目标}

zeroTrace 的目标集中在运行中观测：目标程序继续运行，热点函数的观测成本压低到可接受范围。项目把复杂操作前移到安装阶段：符号解析、远程注入、trampoline 构造、线程组停止和入口 patch 都由 tracer 完成；函数命中后，目标进程内的执行路径集中在现场保存、事件写入和返回链跳转。

这个目标落到实现上，主要对应四个方面：函数命中后的处理留在目标进程内，减少热点函数反复切到内核的成本；参数、返回值和调用现场按目标 ABI 采集，函数本身仍按原来的语义执行；probe 可以在进程运行期间安装、禁用、卸载和更新；多线程场景下，patch 前后先收敛线程组，等所有线程离开入口改写区域后再动手。

\begin{figure}[H]
\centering
\resizebox{0.92\textwidth}{!}{%
\begin{tikzpicture}[
  goal/.style={draw=ztblue, rounded corners=2pt, thick, fill=ztblue!8, minimum width=3.25cm, minimum height=0.82cm, align=center},
  impl/.style={draw=ztgreen, rounded corners=2pt, thick, fill=ztgreen!8, minimum width=3.25cm, minimum height=0.92cm, align=center},
  line/.style={->, thick}
]
  \node[goal] (g1) at (-5.7,1.25) {低命中开销};
  \node[impl] (i1) at (-5.7,0) {用户态 trampoline\\共享 trace buffer};

  \node[goal] (g2) at (-1.9,1.25) {语义保持};
  \node[impl] (i2) at (-1.9,0) {保存现场\\维护返回链};

  \node[goal] (g3) at (1.9,1.25) {动态管理};
  \node[impl] (i3) at (1.9,0) {enable / disable\\untrace / update};

  \node[goal] (g4) at (5.7,1.25) {多线程安全};
  \node[impl] (i4) at (5.7,0) {线程组收敛\\PC 窗口检查};

  \draw[line] (g1) -- (i1);
  \draw[line] (g2) -- (i2);
  \draw[line] (g3) -- (i3);
  \draw[line] (g4) -- (i4);
\end{tikzpicture}
}
\caption{zeroTrace 的主要设计目标}
\label{fig:design-goals}
\end{figure}

\section{需求对应关系与实现特点}

\subsection{赛题要求对应关系}

表 \ref{tab:basic-req} 和表 \ref{tab:advanced-req} 汇总了 zeroTrace 对基础功能 F1--F7 与进阶功能 A1--A5 的支持情况。

\begin{longtable}{@{}P{1.0cm}P{4.2cm}P{9.95cm}@{}}
\caption{基础功能支持情况}\label{tab:basic-req}\\
\toprule
\textbf{编号} & \textbf{赛题要求} & \textbf{zeroTrace 实现}\\
\midrule
\endfirsthead
\toprule
\textbf{编号} & \textbf{赛题要求} & \textbf{zeroTrace 实现}\\
\midrule
\endhead
F1 & 动态库函数探针插入，函数入口命中后直接跳到用户态 trampoline &
支持 \code{read/write/printf} 等动态库函数。x86\_64 使用 \code{jmp qword ptr [rip+0]} 模板，aarch64 使用 \code{ldr x16; br x16} 模板，命中后进入目标进程内的 trampoline。\\
F2 & 函数参数获取，至少前 6 个参数，遵循 ABI &
x86\_64 映射 \code{rdi/rsi/rdx/rcx/r8/r9}，aarch64 映射 \code{x0..x5}；浮点参数映射到统一 \code{fp\_args[]} 槽位。\\
F3 & 函数返回值捕获 &
entry stub 保存真实返回地址并改写返回链；函数返回时先进入 exit stub，记录返回值后再跳回原返回地址。\\
F4 & 探针清理，恢复原始指令并释放远程资源 &
\code{untrace/detach} 按 probe 元数据恢复入口字节；trampoline slot 释放后进入空闲表，pool 全空时远程 \code{munmap}。清理用例会记录入口字节和 \code{/proc/<pid>/maps} 变化。\\
F5 & 探针动态开关 &
\code{disable} 恢复入口 patch，probe 元数据与 slot 继续可用；\code{enable} 复用原 slot 重装 patch。多线程用例会在目标运行期间反复切换。\\
F6 & 多探针共存，同一进程至少 16 个并发探针 &
probe 表容量为 32，自动化测试在同一进程中安装 \code{probe\_fn01..probe\_fn16} 并校验全部命中。\\
F7 & 多线程安全，不导致崩溃或死锁 &
线程组级 \code{seize/interrupt/continue/detach}，运行中刷新新线程；patch 前检查各 TID 的 PC，位于 patch 区间时 single-step 离开。\\
\bottomrule
\end{longtable}

\begin{longtable}{@{}P{1.0cm}P{4.2cm}P{9.95cm}@{}}
\caption{进阶功能支持情况}\label{tab:advanced-req}\\
\toprule
\textbf{编号} & \textbf{进阶要求} & \textbf{zeroTrace 实现}\\
\midrule
\endfirsthead
\toprule
\textbf{编号} & \textbf{进阶要求} & \textbf{zeroTrace 实现}\\
\midrule
\endhead
A1 & x86\_64 + ARM64/aarch64 多架构 &
公共控制层复用，ISA 相关 patch、远程执行、trampoline builder 和 stub 分别落在 \code{src/isa/x86\_64} 与 \code{src/isa/aarch64}。\\
A2 & 条件探针 &
\code{if} 后字符串按完整布尔表达式解析，支持 \code{arg0..arg5}、十进制/十六进制常量、比较、算术、括号和短路求值。\\
A3 & 与 perf/ftrace 事件流合流 &
日志包含 comm、pid/tid、CPU、单调时钟时间戳和事件名；\code{scripts/merge\_trace\_events.py} 能把 zeroTrace、perf script 和 ftrace 风格事件按时间排序。\\
A4 & 探针内调用目标进程函数 &
\code{update <probe> call <callee> ...} 会把 callee 地址和参数来源写入远程 action 表；payload 在目标进程上下文内调用 callee 并写 call event。\\
A5 & 探针热更新 &
filter、call action、enable/disable 状态可在目标运行中更新；热更新用例包含 staged 更新和 live call action 切换。\\
\bottomrule
\end{longtable}

\subsection{同类方案对比}

功能挑战赛道要求说明项目测试结果的功能、性能和创新性，并与类似项目进行对比。zeroTrace 面向“用户态函数命中后的低开销观测”这一场景设计。表 \ref{tab:related-work} 从触发机制、参数表达能力和使用场景三个角度对比常见方案。

\begin{longtable}{@{}P{2.45cm}P{2.9cm}P{3.75cm}P{4.55cm}@{}}
\caption{同类动态追踪方案对比}\label{tab:related-work}\\
\toprule
\textbf{方案} & \textbf{典型机制} & \textbf{优势} & \textbf{zeroTrace 的差异}\\
\midrule
\endfirsthead
\toprule
\textbf{方案} & \textbf{典型机制} & \textbf{优势} & \textbf{zeroTrace 的差异}\\
\midrule
\endhead
kernel uprobe / eBPF & 用户态函数命中后进入内核，事件由内核侧处理 & 内核生态成熟，能与 perf、ftrace、BPF map 等工具组合 & 命中路径会发生用户态到内核态切换；zeroTrace 把 entry/return 处理放在目标进程内，事件写入共享 trace buffer，x86\_64 与 aarch64 benchmark 中额外开销均值分别为 174.45 ns/call 和 237.02 ns/call。\\
ltrace & 依赖动态链接和 PLT 调用路径，按函数签名格式化库函数调用 & 使用简单，配置文件适合描述 libc 函数签名 & ltrace 更偏调试输出；zeroTrace 面向运行中 patch、return probe、多线程安全和低开销 benchmark，同时借鉴 ltrace 风格配置完成参数展示。\\
SystemTap / DTrace 类工具 & 通过脚本化探针语言描述事件和动作 & 表达能力强，适合系统级观测与复杂脚本 & zeroTrace 的脚本能力较轻，重点是用户态函数入口和返回链；条件过滤和 call action 固定在本项目场景内，换取更短的热路径。\\
perf / ftrace & 采集内核事件、采样事件或 tracepoint，输出统一时间线 & 系统事件生态完整，适合性能分析和调度分析 & zeroTrace 输出兼容 perf/ftrace 风格字段，并提供合流脚本；它补充的是普通用户态函数参数和返回值。\\
bpftime / 用户态 eBPF 运行时 & 将部分 eBPF 执行放到用户态，减少内核交互 & 保留 eBPF 编程模型，适合复用 BPF 生态 & zeroTrace 采用 trampoline、stub 和 C payload 直接组织执行流；实现边界更窄，但便于针对赛题要求控制 ABI、返回值和线程组 patch。\\
\bottomrule
\end{longtable}

\subsection{实现特点}

zeroTrace 的实现重点放在函数命中后的执行流、返回链、多线程控制和交互使用上。probe 命中后，执行流直接进入目标进程内的 trampoline 和 payload stub，事件记录也在目标进程内完成；tracer 在后台读取共享内存。return probe 与 entry probe 使用同一条返回链，entry 和 return 通过 \code{call\_id} 配对，既能打印返回值，也能继续扩展为耗时统计。多线程场景中，zeroTrace 的控制范围是整个线程组。CLI 采用调试器式交互，用户可以在目标运行时调整过滤条件、call action 和 probe 状态。

\subsection{创新性分析}

zeroTrace 的主要创新集中在用户态命中路径、返回链维护、线程组安全和实验验证四个方面。

\begin{enumerate}[leftmargin=2.4em]
  \item \textbf{用户态命中路径。} 函数入口 patch 到 trampoline，命中后在目标进程内完成现场保存和事件写入，高频函数的 probe 逻辑直接沿用户态路径执行。
  \item \textbf{return probe 返回链。} x86\_64 通过栈上返回地址改写闭合 entry/return 链路，aarch64 通过 LR 链路完成同样语义；TLS shadow stack 负责按线程保存真实返回地址，entry 与 return 使用 \code{call\_id} 配对。
  \item \textbf{线程组级 patch 安全。} trace runner 在 patch 前刷新目标线程表，纳管运行中新建线程，并检查每个 stopped TID 的 PC；处于入口 patch 窗口的线程会先 single-step 到安全位置。
  \item \textbf{可更新 probe 行为。} filter、enable/disable 和 call action 都能在目标运行期间更新。call action 使用非零 commit token 发布远程 action 快照，payload 读取到的是完整的旧配置、新配置或 disabled 状态。
  \item \textbf{面向提交的验证矩阵。} 测试覆盖 F1--F7、A1--A5、trace buffer wrap、目标退出窗口、热更新压力和多线程压力；benchmark 使用 5 轮重复实验输出 mean/min/max/stdev，并给出 kernel uprobe 对照。
\end{enumerate}

\section{系统架构}

\subsection{总体结构}

zeroTrace 分为控制面和数据面。控制面运行在 tracer 进程中，负责 CLI、会话管理、ptrace 控制、远程注入、trampoline 分配、入口 patch、日志读取与格式化。数据面位于目标进程内部，由函数入口 patch、trampoline、payload stub、payload handler 和 trace buffer 组成。二者通过远程配置区、trace buffer 和 probe 元数据连接。源码导读合并在本章中展开：表 \ref{tab:modules} 对应主要源文件职责，后续小节按一次 probe 的生命周期说明这些模块如何协作。

\begin{figure}[H]
\centering
\resizebox{0.98\textwidth}{!}{%
\begin{tikzpicture}[
  box/.style={draw=black, rounded corners=2pt, thick, fill=ztgray, minimum width=2.8cm, minimum height=0.9cm, align=center},
  hot/.style={draw=ztred, rounded corners=2pt, thick, fill=ztred!8, minimum width=2.8cm, minimum height=0.9cm, align=center},
  data/.style={draw=ztgreen, rounded corners=2pt, thick, fill=ztgreen!8, minimum width=2.8cm, minimum height=0.9cm, align=center},
  line/.style={->, thick},
  dashedline/.style={->, thick, dashed}
]
  \node[box]  (cli)      at (-4.8, 1.25) {CLI\\交互入口};
  \node[box]  (runner)   at (-1.6, 1.25) {trace runner\\会话与日志};
  \node[box]  (injector) at ( 1.6, 1.25) {injector\\ptrace 控制};
  \node[box]  (target)   at ( 4.8, 1.25) {目标进程\\线程组};

  \node[hot]  (patch)    at ( 4.8,-1.15) {函数入口\\patch};
  \node[hot]  (tramp)    at ( 1.6,-1.15) {trampoline\\probe slot};
  \node[hot]  (stub)     at (-1.6,-1.15) {payload stub\\保存现场};
  \node[data] (payload)  at (-4.8,-1.15) {payload handler\\写事件};
  \node[data] (buffer)   at (-8.0,-1.15) {trace buffer\\ring slots};
  \node[box]  (log)      at (-11.2,-1.15) {日志输出\\perf/ftrace 风格};

  \draw[line] (cli) -- (runner);
  \draw[line] (runner) -- (injector);
  \draw[line] (injector) -- (target);
  \draw[line] (target) -- (patch);
  \draw[line] (patch) -- (tramp);
  \draw[line] (tramp) -- (stub);
  \draw[line] (stub) -- (payload);
  \draw[line] (payload) -- (buffer);
  \draw[dashedline] (buffer) -- (log);
  \draw[dashedline, rounded corners=4pt] (runner.south) -- (-1.6,0.28) -- (-8.0,0.28) -- (buffer.north);

  \begin{scope}[on background layer]
    \node[draw=ztblue, dashed, rounded corners=4pt, fit=(cli)(runner)(injector), inner sep=0.35cm, label={[ztblue]above:tracer 进程控制面}] {};
    \node[draw=ztorange, dashed, rounded corners=4pt, fit=(target)(patch)(tramp)(stub)(payload)(buffer), inner sep=0.35cm, label={[ztorange]below:目标进程数据面}] {};
  \end{scope}
\end{tikzpicture}
}
\caption{zeroTrace 控制面与数据面}
\label{fig:architecture}
\end{figure}

图 \ref{fig:architecture} 中红色箭头表示函数命中后最频繁执行的部分。命中阶段主要完成时间戳记录、现场保存和事件写入，执行流留在目标进程内；tracer 的轮询读取共享 buffer，不打断目标进程运行。

\subsection{模块划分}

\begin{longtable}{@{}P{3.6cm}P{3.0cm}P{8.55cm}@{}}
\caption{核心模块职责}\label{tab:modules}\\
\toprule
\textbf{模块} & \textbf{层次} & \textbf{职责}\\
\midrule
\endfirsthead
\toprule
\textbf{模块} & \textbf{层次} & \textbf{职责}\\
\midrule
\endhead
\code{zt\_cli} & 交互层 & 解析用户命令，维护 attach 会话，触发 trace、update、disable、untrace 等操作，并在等待输入时以 50 ms 节流读取新增日志。\\
\code{zt\_trace\_runner} & 控制层 & 串联 injector、payload、trampoline pool、filter、sigconf 和日志输出，是 active trace 的状态中心。\\
\code{zt\_injector} & 注入层 & 管理 ptrace、线程组停止/恢复、新线程纳管、远程内存读写、远程 \code{mmap/munmap}、符号解析、入口 patch 和恢复。\\
\begin{tabular}[t]{@{}l@{}}\code{zt\_trampoline\_}\\\code{pool}\end{tabular} & 资源层 & 管理目标进程中的 executable trampoline pool，按 slot 分配和回收。\\
\code{src/isa/<arch>} & 架构后端 & 提供远程执行、入口跳转模板、patch span 计算、trampoline builder 和 stub。\\
\code{zt\_payload} & 数据面 runtime & 注入目标进程后执行，负责保存现场、维护 TLS shadow stack、写 entry/return/call event。\\
\code{zt\_filter} & 表达式层 & 编译和求值条件探针表达式。\\
\code{zt\_sigconf} & 展示层 & 根据 \code{conf/zttrace.conf} 把原始寄存器槽格式化为函数签名形式。\\
\bottomrule
\end{longtable}

\subsection{核心数据结构}

运行态主要由 probe 表、trace buffer 和 active trace 状态组成。probe 表记录“正在追踪哪些函数”，trace buffer 承接函数命中后的事件，active trace 状态则保存 tracer 与目标进程之间共享的远程资源。

\begin{longtable}{@{}P{3.3cm}P{12.25cm}@{}}
\caption{运行态核心数据结构}\label{tab:runtime-state}\\
\toprule
\textbf{结构} & \textbf{主要内容与作用}\\
\midrule
\endfirsthead
\toprule
\textbf{结构} & \textbf{主要内容与作用}\\
\midrule
\endhead
probe 表 & 保存 \code{probe\_id}、符号名、模块路径、远程函数地址、trampoline 地址、原始入口字节、覆盖长度、状态和 filter。状态显式区分为 \code{resolved/prepared/installed/disabled}，便于 enable、disable、untrace 和热更新复用同一套元数据。\\
trace buffer & 保存 \code{probe\_id}、\code{event\_type}、\code{call\_id}、\code{timestamp\_ns}、\code{tid}、\code{cpu\_id}、通用参数槽、浮点参数槽和 \code{committed\_seq}。entry 记录参数，return 记录返回值，call 记录 probe 内调用目标函数的结果。\\
call action 表 & 由 tracer 写入远程配置区，payload 在 entry 命中时读取。表项包含 \code{enabled} commit token、\code{probe\_id}、callee 地址、参数个数和参数来源。\\
active trace 状态 & 维护 session、远程 payload 地址、trace buffer 地址、payload config 地址、trampoline pool 元数据、日志文件句柄、reader 的 \code{last\_seq} 和运行状态。\\
entry cache & tracer 侧按 \code{call\_id} 缓存 entry 信息，用于 return 配对、条件过滤同步吞掉 return，以及 \code{read/recv/fgets} 等函数在返回后读取 buffer 内容。\\
\bottomrule
\end{longtable}

\subsection{生命周期流程}

一次 probe 从安装到清理分三步：准备运行时、安装入口跳转、消费事件。第一条 probe 会额外完成 payload 注入和共享内存初始化；后续 probe 复用同一份 runtime，再分配新的 trampoline slot。

trace runner 会记录本轮创建的远程 runtime 区域，包括 payload path、trace buffer、payload config 和 trampoline pool。删除最后一个 probe 时，清理顺序是先恢复函数入口，再释放 trampoline slot；如果 pool 已经空闲，再远程 \code{munmap}。随后释放 payload path、trace buffer 和 payload config。目标进程继续存活时，tracer 创建的匿名映射也会随之清理。

\begin{figure}[H]
\centering
\resizebox{0.98\textwidth}{!}{%
\begin{tikzpicture}[
  flowstep/.style={draw=black, rounded corners=2pt, thick, fill=ztgray, minimum width=4.4cm, minimum height=0.85cm, align=center},
  decision/.style={draw=black, diamond, aspect=2.2, thick, fill=ztorange!10, align=center, inner sep=1pt},
  line/.style={->, thick}
]
  \node[flowstep] (attach) at (0,0) {attach 目标进程};
  \node[flowstep] (stop) at (0,-1.35) {停止并收敛线程组};
  \node[decision, minimum width=3.6cm, minimum height=1.35cm] (runtime) at (0,-3.0) {runtime\\已初始化?};
  \node[flowstep] (inject) at (-4.0,-4.65) {远程 dlopen payload\\分配 buffer/config};
  \node[flowstep] (reuse) at (4.0,-4.65) {复用已有 runtime};
  \node[flowstep] (symbol) at (0,-6.3) {解析符号地址\\读取原始前导指令};
  \node[flowstep] (build) at (0,-7.65) {分配 trampoline slot\\重定位原始指令};
  \node[flowstep] (patch) at (0,-9.0) {入口 patch 到 trampoline};
  \node[flowstep] (resume) at (0,-10.35) {恢复目标线程组};
  \node[flowstep] (poll) at (0,-11.7) {异步读取 trace buffer};

  \draw[line] (attach) -- (stop);
  \draw[line] (stop) -- (runtime);
  \draw[line] (runtime.south west) -- node[above left]{否} (inject.north);
  \draw[line] (runtime.south east) -- node[above right]{是} (reuse.north);
  \draw[line] (inject) |- (symbol);
  \draw[line] (reuse) |- (symbol);
  \draw[line] (symbol) -- (build);
  \draw[line] (build) -- (patch);
  \draw[line] (patch) -- (resume);
  \draw[line] (resume) -- (poll);
\end{tikzpicture}
}
\caption{probe 安装与运行流程}
\label{fig:probe-lifecycle}
\end{figure}

\begin{figure}[H]
\centering
\resizebox{0.94\textwidth}{!}{%
\begin{tikzpicture}[
  slot/.style={draw=black, thick, minimum width=1.25cm, minimum height=0.75cm, align=center, fill=ztgray},
  used/.style={draw=ztred, thick, minimum width=1.25cm, minimum height=0.75cm, align=center, fill=ztred!10},
  free/.style={draw=ztgreen, thick, minimum width=1.25cm, minimum height=0.75cm, align=center, fill=ztgreen!10},
  box/.style={draw=black, rounded corners=2pt, thick, fill=ztgray, minimum width=3.3cm, minimum height=0.82cm, align=center},
  action/.style={draw=ztblue, rounded corners=2pt, thick, fill=ztblue!8, minimum width=3.3cm, minimum height=0.82cm, align=center},
  line/.style={->, thick}
]
  \node[box] (pool) at (0,1.55) {远程 executable pool\\一次 mmap，按 slot 切分};
  \node[used] (s0) at (-4.55,0) {slot 0\\probe 1};
  \node[used] (s1) at (-3.15,0) {slot 1\\probe 2};
  \node[free] (s2) at (-1.75,0) {slot 2\\free};
  \node[used] (s3) at (-0.35,0) {slot 3\\probe 7};
  \node[free] (s4) at (1.05,0) {slot 4\\free};
  \node[slot] (s5) at (2.45,0) {\ldots};
  \node[free] (s6) at (3.85,0) {slot n\\free};

  \node[action] (alloc) at (-3.2,-1.6) {trace/enable\\free slot 变为 used};
  \node[action] (release) at (0.35,-1.6) {untrace/detach\\used slot 回到 free};
  \node[box] (unmap) at (3.9,-1.6) {全部 slot 空闲\\远程 munmap};

  \draw[line] (pool.south) -- (s3.north);
  \draw[line, ztgreen] (s2.south) -- (alloc.north);
  \draw[line, ztred] (s3.south) -- (release.north);
  \draw[line, ztorange] (release.east) -- (unmap.west);
\end{tikzpicture}
}
\caption{trampoline pool 与 slot 回收方式}
\label{fig:trampoline-pool}
\end{figure}

\section{关键实现}

\subsection{入口 patch 与 trampoline}

zeroTrace 改写函数入口时写入固定长度的跳转模板。x86\_64 使用 14 字节 RIP 相对间接跳转模板，模板内容为 \code{ff 25 00 00 00 00} 后接 64 位 trampoline 地址；aarch64 使用 16 字节 \code{ldr/br} 模板，并按 AAPCS64 选择 \code{x16/ip0} 作为跳转寄存器。两种模板都会把执行流送入远程 trampoline，再由 trampoline 负责执行被搬运的前导指令并跳回原函数。

\begin{lstlisting}
x86_64 entry patch:
    jmp qword ptr [rip + 0]
    .quad trampoline_addr

aarch64 entry patch:
    ldr x16, #8
    br  x16
    .quad trampoline_addr
\end{lstlisting}

入口 patch 会覆盖目标函数前导若干条指令。x86\_64 后端使用 Capstone 按指令边界累计长度，直到满足模板长度；aarch64 指令固定 4 字节，按 4 字节对齐累计。被覆盖的原始指令会复制到 trampoline 中执行，之后再跳回 \code{func + orig\_len}。

\begin{figure}[H]
\centering
\resizebox{0.95\textwidth}{!}{%
\begin{tikzpicture}[
  box/.style={draw=black, rounded corners=2pt, thick, fill=ztgray, minimum width=3.2cm, minimum height=0.85cm, align=center},
  hot/.style={draw=ztred, rounded corners=2pt, thick, fill=ztred!8, minimum width=3.2cm, minimum height=0.85cm, align=center},
  line/.style={->, thick}
]
  \node[box] (entry) at (-4.2,0) {原函数入口\\patch};
  \node[hot] (tramp) at (0,0) {trampoline\\push probe id};
  \node[hot] (entry_stub) at (4.2,0) {entry stub\\采集参数};
  \node[box] (orig) at (4.2,-1.75) {重定位后的\\原始前导指令};
  \node[box] (body) at (0,-1.75) {原函数主体\\func + orig\_len};
  \node[hot] (exit_stub) at (0,-3.5) {exit stub\\采集返回值};
  \node[box] (ret) at (4.2,-3.5) {真实返回地址};

  \draw[line] (entry) -- (tramp);
  \draw[line] (tramp) -- (entry_stub);
  \draw[line] (entry_stub) -- (orig);
  \draw[line] (orig) -- (body);
  \draw[line] (body) -- (exit_stub);
  \draw[line] (exit_stub) -- (ret);
\end{tikzpicture}
}
\caption{函数命中后的控制流}
\label{fig:runtime-flow}
\end{figure}

\subsection{前导指令重定位}

trampoline builder 按原函数前导语义处理指令。普通指令可以直接搬运；依赖当前位置的指令需要重定位或改写。

\begin{longtable}{@{}P{3.3cm}P{12.25cm}@{}}
\caption{前导指令重定位策略}\label{tab:relocation}\\
\toprule
\textbf{指令类型} & \textbf{处理方式}\\
\midrule
\endfirsthead
\toprule
\textbf{指令类型} & \textbf{处理方式}\\
\midrule
\endhead
普通指令 & 与当前 PC 无关，可原样复制。\\
x86\_64 RIP 相对内存访问 & 计算 \code{old\_target = old\_next\_ip + old\_disp}，再回填 \code{new\_disp = old\_target - new\_next\_ip}。\\
x86\_64 \code{call rel32} & 改写为绝对调用模板，消除相对偏移距离限制。\\
x86\_64 \code{jmp rel8/rel32} & 改写为绝对跳转模板。\\
x86\_64 \code{jcc rel8/rel32} & 展开为“条件反转 + 跳过绝对跳转块”的模板。\\
aarch64 \code{adr/adrp} & 根据 trampoline 地址重新编码目标页或目标地址。\\
aarch64 分支类指令 & 对 \code{b/bl/b.cond/cbz/cbnz/tbz/tbnz} 进行重定位；遇到暂不支持的模式时返回错误，不生成 trampoline。\\
aarch64 literal load & 这类指令的目标地址依赖 PC；安装阶段识别后停止生成 trampoline，入口 patch 不会写入目标进程。\\
\bottomrule
\end{longtable}

\subsection{return probe 与返回链接管}

return probe 的关键约束是保持被测函数自己的栈语义。zeroTrace 在 entry stub 中保存真实返回地址，再把函数的返回目标改成 exit stub。函数主体照常执行，ret 落到 exit stub 后记录返回寄存器，再从 TLS shadow stack 中取回真实返回地址。

x86\_64 上，调用者的返回地址位于栈上，entry stub 可以直接改写对应槽位。aarch64 上，返回地址通常在 \code{x30/lr} 中，trampoline 会保存原 LR，entry stub 将 exit stub 写入返回链，真实 LR 同样进入线程本地 shadow stack。两套实现细节不同，但上层 payload 看到的事件模型一致。

x86\_64 的 exit stub 入口来自原函数 \code{ret}，不同于普通 \code{call}。为了让 C handler 仍能按同一份上下文结构读取返回值，它会先补出返回地址槽和占位槽，再执行通用寄存器保存序列。记录 return event 后，exit stub 从 TLS shadow stack 弹出真实返回地址，写回预留槽位，最后通过 \code{ret} 回到原调用者。aarch64 的处理放在 \code{x15/x30} 这条 LR 链路上：\code{x15} 保存真实 LR，entry stub 返回前把 exit stub 地址带回 trampoline，trampoline 再把它写入 \code{x30}，原函数返回时进入 exit stub，最后恢复真实 LR。

\begin{figure}[H]
\centering
\resizebox{0.95\textwidth}{!}{%
\begin{tikzpicture}[
  cell/.style={draw=black, minimum width=7.8cm, minimum height=0.55cm, align=left, fill=ztgray},
  mark/.style={draw=ztred, very thick, fill=ztred!8},
  note/.style={align=left, font=\small, inner sep=1pt}
]
  \node[cell] (r0) {高地址：真实调用者返回地址};
  \node[cell, below=0pt of r0] (r1) {probe id};
  \node[cell, below=0pt of r1] (r2) {trampoline call entry\_stub 返回地址};
  \node[cell, below=0pt of r2] (r3) {保存的 FLAGS / 通用寄存器};
  \node[cell, below=0pt of r3] (r4) {紧凑 XMM / FP 快照};
  \node[cell, below=0pt of r4] (r5) {低地址：stub 临时栈帧};

  \node[note, right=0.9cm of r0] (n0) {保存到 TLS shadow stack};
  \node[note, right=0.9cm of r1] (n1) {查找 probe 元数据};
  \node[note, right=0.9cm of r4] (n2) {记录用于展示的 XMM/FP 槽};

  \draw[->, thick, ztred] (r0.east) -- (n0.west);
  \draw[->, thick, ztblue] (r1.east) -- (n1.west);
  \draw[->, thick, ztgreen] (r4.east) -- (n2.west);
\end{tikzpicture}
}
\caption{x86\_64 entry stub 的逻辑栈布局}
\label{fig:stack-layout}
\end{figure}

\begin{figure}[H]
\centering
\resizebox{0.98\textwidth}{!}{%
\begin{tikzpicture}[
  nodebox/.style={draw=black, rounded corners=2pt, thick, fill=ztgray, minimum width=3.15cm, minimum height=0.8cm, align=center, font=\small},
  hotbox/.style={draw=ztred, rounded corners=2pt, thick, fill=ztred!8, minimum width=3.15cm, minimum height=0.8cm, align=center, font=\small},
  storebox/.style={draw=ztgreen, rounded corners=2pt, thick, fill=ztgreen!8, minimum width=3.15cm, minimum height=0.8cm, align=center, font=\small},
  line/.style={->, thick}
]
  \node[nodebox] (caller) at (-5.1,1.55) {调用者\\LR 写入 x30};
  \node[hotbox] (tramp) at (-1.7,1.55) {trampoline\\x15 = 原 LR};
  \node[hotbox] (entry) at (1.7,1.55) {entry\_stub\\记录参数};
  \node[storebox] (shadow) at (5.1,1.55) {TLS shadow stack\\保存真实 LR};

  \node[hotbox] (setexit) at (1.7,0.0) {entry\_stub\\x15 = exit\_stub};
  \node[hotbox] (setlr) at (-1.7,0.0) {trampoline\\x30 = x15};
  \node[nodebox] (body) at (-5.1,0.0) {原函数主体\\执行并返回};

  \node[hotbox] (exit) at (-1.7,-1.55) {exit\_stub\\记录返回值};
  \node[storebox] (restore) at (1.7,-1.55) {弹出真实 LR\\写回 x30};
  \node[nodebox] (ret) at (5.1,-1.55) {返回调用者\\ret};

  \draw[line] (caller.east) -- (tramp.west);
  \draw[line] (tramp.east) -- (entry.west);
  \draw[line] (entry.east) -- (shadow.west);
  \draw[line] (entry.south) -- (setexit.north);
  \draw[line] (setexit.west) -- (setlr.east);
  \draw[line] (setlr.west) -- (body.east);
  \draw[line] (body.south east) -- (exit.north west);
  \draw[line] (exit.east) -- (restore.west);
  \draw[line] (restore.east) -- (ret.west);
\end{tikzpicture}
}
\caption{aarch64 return probe 的 LR 接管流程}
\label{fig:aarch64-lr-chain}
\end{figure}

浮点上下文按 ABI 参数和返回值槽位维护。x86\_64 stub 对应 \code{xmm0..xmm7}，aarch64 stub 对应 FP/SIMD 参数槽位，并保留 \code{fpcr/fpsr/nzcv} 等状态字段。payload 读取统一的 \code{fp\_args[]} 槽位，上层格式化逻辑只处理统一后的参数槽，使两套 ISA 后端保持一致的事件模型。

\subsection{trace buffer 发布协议}

payload 与 tracer 之间通过目标进程内的 ring buffer 传递事件。payload 写入时先用原子序号 reserve 一个 slot，直接在 slot 内填充 event 字段，最后写入 \code{committed\_seq}。tracer 读取时先读 header；序号有变化时再按范围读取事件区，并校验 slot 前后的提交序号。

\begin{figure}[H]
\centering
\resizebox{0.88\textwidth}{!}{%
\begin{tikzpicture}[
  box/.style={draw=black, rounded corners=2pt, thick, fill=ztgray, minimum width=3.4cm, minimum height=0.85cm, align=center},
  data/.style={draw=ztgreen, rounded corners=2pt, thick, fill=ztgreen!8, minimum width=3.4cm, minimum height=0.85cm, align=center},
  warn/.style={draw=ztorange, rounded corners=2pt, thick, fill=ztorange!10, minimum width=3.4cm, minimum height=0.85cm, align=center},
  line/.style={->, thick}
]
  \node[box]  (reserve) at (-4.0, 1.4) {reserve seq};
  \node[data] (fill)    at ( 0.0, 1.4) {直接写 event slot};
  \node[data] (commit)  at ( 4.0, 1.4) {commit seq};
  \node[box]  (header)  at ( 4.0, 0.0) {reader 读取 header};
  \node[box]  (range)   at ( 0.0, 0.0) {计算可见序列范围};
  \node[warn] (verify)  at (-4.0, 0.0) {校验 committed\_seq};
  \node[data] (consume) at (-4.0,-1.4) {格式化并写日志};
  \node[warn] (lost)    at ( 0.0,-1.4) {落后过多时\\输出 lost record};

  \draw[line] (reserve) -- (fill);
  \draw[line] (fill) -- (commit);
  \draw[line] (commit) -- (header);
  \draw[line] (header) -- (range);
  \draw[line] (range) -- (verify);
  \draw[line] (verify) -- (consume);
  \draw[line] (range) -- (lost);
\end{tikzpicture}
}
\caption{trace buffer 写入与消费协议}
\label{fig:ring-buffer}
\end{figure}

这套协议主要处理两种异常情况：slot 尚未提交时，reader 停在对应序列，下一轮再读，不消费半写入事件；reader 落后超过 ring capacity 时，日志中会出现 \code{ztrace:lost}，并写出丢失范围。

tracer 读取目标内存时优先使用 \code{process\_vm\_readv}，失败后再回退到 ptrace 逐字读取。正常轮询先读取 header 中的 magic 和写入序号；序号未变化时停在 header。发生 wrap 时，reader 会跳到仍可见的最早序列，并在日志中记录丢失范围。目标进程退出窗口需要单独处理：如果远程读取失败，trace runner 会结合 \code{kill(pid, 0)}、\code{/proc/<pid>/stat} 中的 zombie/dead 状态、trace buffer 映射是否仍出现在 \code{/proc/<pid>/maps} 中以及短暂确认窗口判断是否属于正常退出。确认退出后关闭 active trace，后续重复 poll 返回完成态，将正常退出和 trace 错误区分开。

\subsection{条件过滤与热更新}

条件过滤放在 tracer 消费 entry event 时求值，目标进程里的 payload 继续负责采集快照和写事件。若 entry 不满足条件，trace runner 会把该 \code{call\_id} 标记为 suppressed；后续 return event 也按同一个标记丢弃。

\begin{figure}[H]
\centering
\resizebox{0.96\textwidth}{!}{%
\begin{tikzpicture}[
  box/.style={draw=black, rounded corners=2pt, thick, fill=ztgray, minimum width=2.85cm, minimum height=0.78cm, align=center, font=\small},
  data/.style={draw=ztgreen, rounded corners=2pt, thick, fill=ztgreen!8, minimum width=2.85cm, minimum height=0.78cm, align=center, font=\small},
  warn/.style={draw=ztorange, rounded corners=2pt, thick, fill=ztorange!10, minimum width=2.85cm, minimum height=0.78cm, align=center, font=\small},
  decision/.style={draw=black, diamond, aspect=2.55, thick, fill=ztorange!10, align=center, inner sep=1pt, font=\small},
  line/.style={->, thick}
]
  \node[data] (entry) at (-6.4,1.60) {entry event\\参数快照};
  \node[box] (cache) at (-3.1,1.60) {entry cache\\保存 call\_id};
  \node[decision, minimum width=2.75cm, minimum height=1.05cm] (filter) at (0.2,1.60) {filter\\通过?};
  \node[data] (elog) at (3.65,1.60) {写 entry 日志};
  \node[warn] (sup) at (0.2,0.15) {suppressed set\\记录 call\_id};
  \node[warn] (edrop) at (3.65,0.15) {丢弃 entry};

  \node[data] (ret) at (-6.4,-3.15) {return event\\返回值快照};
  \node[box] (lookup) at (-3.1,-3.15) {按 call\_id\\查 entry 状态};
  \node[decision, minimum width=2.75cm, minimum height=1.05cm] (keep) at (0.2,-3.15) {entry\\已通过?};
  \node[data] (rlog) at (3.65,-3.15) {写 return 日志};
  \node[warn] (rdrop) at (3.65,-4.55) {丢弃 return};

  \draw[line] (entry.east) -- (cache.west);
  \draw[line] (cache.east) -- (filter.west);
  \draw[line] (filter.east) -- node[above]{是} (elog.west);
  \draw[line] (filter.south) -- node[right]{否} (sup.north);
  \draw[line] (sup.east) -- (edrop.west);
  \draw[line] (ret.east) -- (lookup.west);
  \draw[line] (lookup.east) -- (keep.west);
  \draw[line] (keep.east) -- node[above]{是} (rlog.west);
  \draw[line] (keep.south east) |- node[pos=0.28,right]{否} (rdrop.west);
\end{tikzpicture}
}
\caption{条件过滤与 return 事件配对}
\label{fig:filter-flow}
\end{figure}

\begin{lstlisting}
trace write if arg0 == 1 && arg2 > 0
trace probe_fn01 if arg0 >= 0x10 && arg0 < 0x20
update probe_fn01 if arg0 >= 100
update probe_fn01 clear
\end{lstlisting}

\code{update} 采用替换语义。filter 更新发生在 tracer 侧，目标进程继续运行；call action 更新会写入远程 action 表，并通过非零 commit token 发布完整快照；enable/disable 涉及入口 patch，会短暂停住线程组。

filter 解析借鉴 NEMU 表达式求值思路：先把 \code{if} 后面的字符串 token 化，再按优先级递归求值。表达式支持 \code{arg0..arg5}、十进制和十六进制常量、比较、算术、括号、\code{\&\&}/\code{||}/\code{!}。运行期按 entry event 中的参数槽求值，\code{\&\&} 和 \code{||} 使用短路语义；例如 \code{arg0 == 0 || 10 / arg0 > 1} 左侧成立后直接返回 true，右侧除法被跳过。

call action 更新需要同步目标进程中的 action 表。tracer 写入时先把对应 slot 的 \code{enabled} token 清零，再写入 \code{probe\_id}、callee 地址和参数来源，最后用新的非零 token 发布。payload 读取时先后两次读取 token，中间拷贝表项快照；只有两次 token 相同且非零时才接受该快照。这样运行中切换 callee 或清空 action 时，payload 看到的是旧配置、新配置或 disabled。action 表按实际 \code{probe\_id} 查找匹配项，长期运行中即使出现 \code{probe\_id \% 32} 碰撞，也能区分不同 probe。

\subsection{签名配置与参数解码}

zeroTrace 参考 \code{ltrace.conf} 的思路，把常见函数签名整理成简化配置。用户补一行签名即可获得参数名、类型和返回值格式化；未配置函数则退回寄存器风格显示。

\begin{lstlisting}
function_name(arg_type arg_name, arg_type arg_name, ...) -> return_type

puts(const char *s) -> int
read(int fd, buffer buf, size_t count) -> long
write(int fd, const buffer buf, size_t count) -> long
malloc(size_t size) -> void *
fp_mix(double a, double b) -> double
\end{lstlisting}

参数解码按 ABI 参数槽位进行。整数和指针参数来自通用参数槽，浮点参数来自浮点参数槽；返回值也按整数返回槽和浮点返回槽分开读取。

字符串参数会通过远程内存读取转义为可打印内容，buffer 类型会结合相邻长度参数或返回值输出部分内容。可变参数函数中，参数名为 \code{fmt} 的格式化字符串会触发 varargs 展开，格式化逻辑根据签名参数名判断。实现范围覆盖寄存器快照内的整数、指针和浮点可变参数；已经落到调用栈上的 varargs 需要额外的栈快照机制，本版先聚焦寄存器快照内的参数。

\begin{table}[H]
\centering
\caption{签名配置类型与输出行为}
\label{tab:sigconf}
\begin{tabularx}{\textwidth}{P{3.2cm}Y}
\toprule
\textbf{类型/约定} & \textbf{行为}\\
\midrule
\code{int/long/size\_t} & 从通用参数槽读取并按数值输出。\\
\code{void * / addr} & 从通用参数槽读取并按地址输出。\\
\code{char * / string} & 读取目标进程远程字符串，非打印字符输出为 \code{\textbackslash xNN}。\\
\code{buffer} & 结合长度参数或返回值展示读写内容片段。\\
\code{float/double} & 从浮点参数槽或浮点返回槽解码。\\
\code{fmt} 参数名 & 读取 format string，并展开寄存器内 varargs。\\
\bottomrule
\end{tabularx}
\end{table}

\subsection{多线程 patch 安全}

多线程安全的关键在于避开 patch 与线程执行窗口互相踩踏。zeroTrace 的线程组控制分为 attach、运行态刷新和 patch 前收敛三部分。

\begin{enumerate}[leftmargin=2.4em]
  \item attach 时枚举 \code{/proc/<pid>/task}，对当前线程执行 \code{PTRACE\_SEIZE}。
  \item interrupt 阶段多轮刷新线程表，把新创建线程纳入控制，直到线程表稳定且所有 tracked TID 均 stopped。
  \item 运行期间 \code{zt\_trace\_poll()} 继续消费 ptrace event，并把普通 signal-delivery-stop 的信号透传回目标线程。
  \item 改写或恢复入口 patch 前，读取每个 stopped TID 的 PC；若 PC 位于任一 patch 区间，则 single-step 到安全位置。
  \item patch 完成后恢复线程组，或在 detach 时逐个解除 ptrace 控制。
\end{enumerate}

这套流程专门处理运行中创建线程的窗口。每轮 interrupt 后都会重新读取目标线程表，发现新 TID 就继续 seize 和 interrupt；线程表稳定、tracked TID 全部停止后，才进入 patch 阶段。single-step 避让 patch 区间时会记录该线程原本需要透传的 signal，离开 patch 区间后继续交还给目标程序。

\begin{figure}[H]
\centering
\resizebox{0.98\textwidth}{!}{%
\begin{tikzpicture}[
  box/.style={draw=black, rounded corners=2pt, thick, fill=ztgray, minimum width=2.55cm, minimum height=0.70cm, align=center, font=\small},
  hot/.style={draw=ztred, rounded corners=2pt, thick, fill=ztred!8, minimum width=2.55cm, minimum height=0.70cm, align=center, font=\small},
  data/.style={draw=ztgreen, rounded corners=2pt, thick, fill=ztgreen!8, minimum width=2.55cm, minimum height=0.70cm, align=center, font=\small},
  line/.style={->, thick},
  sep/.style={ztline, dashed, thick}
]
  \path[use as bounding box] (-6.55,-2.15) rectangle (6.65,2.25);

  \draw[sep] (-5.55,0.95) -- (6.35,0.95);
  \draw[sep] (-5.55,-0.45) -- (6.35,-0.45);

  \node[font=\small\bfseries, anchor=east] at (-5.75,1.65) {tracer};
  \node[font=\small\bfseries, anchor=east] at (-5.75,0.25) {线程组};
  \node[font=\small\bfseries, anchor=east] at (-5.75,-1.15) {patch 区间};

  \node[box] (seize) at (-4.15,1.65) {枚举 task\\PTRACE\_SEIZE};
  \node[box] (interrupt) at (-1.15,1.65) {interrupt\\刷新线程表};
  \node[hot] (pc) at (1.85,1.65) {检查每个 TID\\当前 PC};
  \node[data] (resume) at (4.85,1.65) {continue\\恢复运行};

  \node[box] (running) at (-4.15,0.25) {running};
  \node[hot] (stopped) at (-1.15,0.25) {all stopped};
  \node[hot] (step) at (1.85,0.25) {命中 patch 窗口\\single-step};
  \node[box] (run2) at (4.85,0.25) {running};

  \node[hot] (old) at (-1.15,-1.15) {旧入口字节};
  \node[data] (new) at (1.85,-1.15) {写入/恢复\\入口 patch};
  \node[data] (safe) at (4.85,-1.15) {安全窗口结束};

  \draw[line] (seize.east) -- (interrupt.west);
  \draw[line] (interrupt.east) -- (pc.west);
  \draw[line] (pc.east) -- (resume.west);
  \draw[line] (running.east) -- (stopped.west);
  \draw[line] (stopped.east) -- (step.west);
  \draw[line] (step.east) -- (run2.west);
  \draw[line] (old.east) -- (new.west);
  \draw[line] (new.east) -- (safe.west);
\end{tikzpicture}
}
\caption{线程组收敛与 patch 安全窗口}
\label{fig:thread-patch-safety}
\end{figure}

因此，stop、continue、enable、disable、untrace 都会覆盖整个线程组。自动化测试记录多线程 entry/return 配平、动态开关次数、线程组停止窗口和恢复后的心跳进度。

\subsection{多架构后端}

ptrace 注入的高层流程与 CPU 架构关系不大，但远程 syscall、远程函数调用、入口 patch 模板、trampoline 重定位和 stub 现场保存都强依赖 ISA。zeroTrace 因此把架构相关实现收敛到 \code{src/isa/<arch>}。

\begin{figure}[H]
\centering
\resizebox{0.92\textwidth}{!}{%
\begin{tikzpicture}[
  box/.style={draw=black, rounded corners=2pt, thick, fill=ztgray, minimum width=3.8cm, minimum height=0.82cm, align=center},
  api/.style={draw=ztblue, rounded corners=2pt, thick, fill=ztblue!8, minimum width=4.1cm, minimum height=0.82cm, align=center},
  arch/.style={draw=ztgreen, rounded corners=2pt, thick, fill=ztgreen!8, minimum width=3.7cm, minimum height=0.82cm, align=center},
  hot/.style={draw=ztred, rounded corners=2pt, thick, fill=ztred!8, minimum width=3.7cm, minimum height=0.82cm, align=center},
  line/.style={->, thick}
]
  \node[box] (runner) at (0,2.4) {公共控制层\\trace runner / injector};
  \node[api] (ops) at (0,0.95) {架构接口\\patch / call / syscall / relocate};
  \node[arch] (x86) at (-3.9,-0.7) {x86\_64 后端\\14B patch，RIP 重定位};
  \node[arch] (arm) at (3.9,-0.7) {aarch64 后端\\16B patch，PC 重定位};
  \node[hot] (xstub) at (-3.9,-2.1) {stub.S\\SysV ABI 现场};
  \node[hot] (astub) at (3.9,-2.1) {stub.S\\AAPCS64 现场};

  \draw[line] (runner.south) -- (ops.north);
  \draw[line] (ops.south west) -- (x86.north);
  \draw[line] (ops.south east) -- (arm.north);
  \draw[line] (x86.south) -- (xstub.north);
  \draw[line] (arm.south) -- (astub.north);
\end{tikzpicture}
}
\caption{公共控制层与 ISA 后端边界}
\label{fig:arch-backends}
\end{figure}

\begin{table}[H]
\centering
\caption{架构后端分工}
\label{tab:arch-boundary}
\begin{tabularx}{\textwidth}{P{4.5cm}Y}
\toprule
\textbf{接口/文件} & \textbf{职责}\\
\midrule
\code{arch.c} & 远程 syscall、远程函数调用、patch 长度、patch span 计算和入口跳转安装。\\
\code{trampoline\_manager.c} & 构造该 ISA 下的 trampoline，并处理前导指令重定位。\\
\code{stub.S} & 保存/恢复现场，调用 payload handler，维护 return probe 返回链。\\
\begin{tabular}[t]{@{}l@{}}\code{isa/common/}\\\code{zt\_remote\_exec.c}\end{tabular} & 提供远程执行的通用控制流：保存现场、写入临时 stub、继续目标进程、等待陷入点、读取返回值并恢复现场；各 ISA 后端补充寄存器布局和 stub 编码。\\
\begin{tabular}[t]{@{}l@{}}\code{scripts/}\\\code{check\_arch\_config.py}\end{tabular} & 检查 Makefile 在 \code{ARCH=x86\_64/aarch64} 下是否选中正确后端和测试集合。\\
\bottomrule
\end{tabularx}
\end{table}

\begin{table}[H]
\centering
\caption{x86\_64 与 aarch64 后端实现对照}
\label{tab:arch-compare}
\begin{tabularx}{\textwidth}{P{3.1cm}Y Y}
\toprule
\textbf{项目} & \textbf{x86\_64} & \textbf{aarch64}\\
\midrule
入口 patch & 14 字节模板，\code{jmp qword ptr [rip+0]} 后接 64 位 trampoline 地址。 & 16 字节模板，\code{ldr x16, \#8; br x16} 后接 64 位 trampoline 地址。\\
参数槽 & SysV ABI 前 6 个整数参数来自 \code{rdi/rsi/rdx/rcx/r8/r9}。 & AAPCS64 前 6 个整数参数来自 \code{x0..x5}。\\
返回值槽 & 整数返回值来自 \code{rax}，浮点返回值来自 \code{xmm0}。 & 整数返回值来自 \code{x0}，浮点返回值来自 \code{q0/d0} 对应槽位。\\
return probe & entry stub 改写栈上的返回地址，原返回地址进入 TLS shadow stack。 & trampoline 保存原 \code{x30/lr}，entry stub 通过 \code{x15} 把 \code{exit\_stub} 地址带回并写入 \code{x30}。\\
远程执行 & 使用 \code{PTRACE\_GETREGS/SETREGS} 操作 \code{user\_regs\_struct}。 & 使用 \code{PTRACE\_GETREGSET/SETREGSET} 操作 \code{struct user\_pt\_regs}。\\
临时调用 stub & 设置参数寄存器后执行目标函数，返回后用断点停回 tracer。 & \code{x0/x1} 传参，\code{blr x16} 调用目标函数，\code{brk \#0} 停回 tracer。\\
前导指令处理 & 需要处理可变长指令、RIP 相对寻址、相对 call/jmp/jcc。 & 指令固定 4 字节，重点处理 PC 相对地址和分支立即数。\\
\bottomrule
\end{tabularx}
\end{table}

\begin{table}[H]
\centering
\caption{架构特性对 zeroTrace 实现的影响}
\label{tab:arch-design-impact}
\begin{tabularx}{\textwidth}{P{3.0cm}Y Y}
\toprule
\textbf{架构特性} & \textbf{x86\_64 上的实现影响} & \textbf{aarch64 上的实现影响}\\
\midrule
指令长度 & 指令长度可变，入口 patch 前需要反汇编并累计完整指令边界。 & 指令固定 4 字节，patch span 按 4 字节对齐累计，边界判断更直接。\\
绝对跳转形式 & 可用 RIP 相对间接跳转，patch 模板短，但 trampoline 内要处理 RIP 相对指令。 & 使用 literal load 加 \code{br x16} 完成绝对跳转，模板固定为 16 字节。\\
返回地址位置 & 普通函数返回地址在调用栈上，entry stub 改写该栈槽即可接管 return probe。 & 返回地址在 \code{x30/lr} 中，trampoline 和 stub 需要显式传递、保存和替换 LR。\\
系统调用约定 & 系统调用号和参数寄存器与 SysV 函数 ABI 不同，远程 syscall stub 需要单独设置。 & \code{x8} 保存 syscall number，\code{x0..x5} 保存参数，返回值仍在 \code{x0}。\\
浮点参数槽 & SysV ABI 使用 \code{xmm0..xmm7} 传递常见浮点参数。 & AAPCS64 使用 \code{v0..v7/q0..q7} 传递常见浮点参数。\\
条件码 & \code{rflags} 同时承载算术状态和条件跳转依据。 & \code{nzcv} 承载条件码，stub 需要配合 \code{mrs/msr} 保存恢复。\\
\bottomrule
\end{tabularx}
\end{table}

\begin{figure}[H]
\centering
\resizebox{0.88\textwidth}{!}{%
\begin{tikzpicture}[
  head/.style={draw=ztblue, rounded corners=2pt, thick, fill=ztblue!8, minimum width=4.25cm, minimum height=0.72cm, align=center, font=\small\bfseries},
  flowbox/.style={draw=black, rounded corners=2pt, thick, fill=ztgray, minimum width=4.25cm, minimum height=0.72cm, align=center, font=\small},
  hot/.style={draw=ztred, rounded corners=2pt, thick, fill=ztred!8, minimum width=4.25cm, minimum height=0.72cm, align=center, font=\small},
  tail/.style={draw=ztgreen, rounded corners=2pt, thick, fill=ztgreen!8, minimum width=4.25cm, minimum height=0.72cm, align=center, font=\small},
  line/.style={->, thick}
]
  \node[head] (xh) at (-3.2,2.95) {x86\_64 命中路径};
  \node[flowbox] (x1) at (-3.2,2.05) {14B entry patch};
  \node[hot]  (x2) at (-3.2,1.15) {trampoline 调用 entry stub};
  \node[hot]  (x3) at (-3.2,0.25) {entry stub 改写栈上返回地址};
  \node[flowbox] (x4) at (-3.2,-0.65) {执行原函数主体};
  \node[hot]  (x5) at (-3.2,-1.55) {ret 进入 exit stub};
  \node[tail] (x6) at (-3.2,-2.45) {恢复真实返回地址};

  \node[head] (ah) at (3.2,2.95) {aarch64 命中路径};
  \node[flowbox] (a1) at (3.2,2.05) {16B entry patch};
  \node[hot]  (a2) at (3.2,1.15) {trampoline 保存原 LR};
  \node[hot]  (a3) at (3.2,0.25) {entry stub 设置 x30};
  \node[flowbox] (a4) at (3.2,-0.65) {执行原函数主体};
  \node[hot]  (a5) at (3.2,-1.55) {ret 进入 exit stub};
  \node[tail] (a6) at (3.2,-2.45) {写回真实 LR};

  \foreach \from/\to in {xh/x1,x1/x2,x2/x3,x3/x4,x4/x5,x5/x6,ah/a1,a1/a2,a2/a3,a3/a4,a4/a5,a5/a6} {
    \draw[line] (\from) -- (\to);
  }
  \draw[dashed, thick, ztline] (0,-2.85) -- (0,3.35);
\end{tikzpicture}
}
\caption{x86\_64 与 aarch64 命中路径对比}
\label{fig:arch-hotpath-compare}
\end{figure}

\begin{figure}[H]
\centering
\resizebox{0.88\textwidth}{!}{%
\begin{tikzpicture}[
  title/.style={draw=ztblue, rounded corners=2pt, thick, fill=ztblue!8, minimum width=4.85cm, minimum height=0.75cm, align=center, font=\small\bfseries},
  cell/.style={draw=black, rounded corners=1pt, thick, fill=ztgray, minimum width=4.85cm, minimum height=0.72cm, align=center, font=\small},
  hot/.style={draw=ztred, rounded corners=1pt, thick, fill=ztred!8, minimum width=4.85cm, minimum height=0.72cm, align=center, font=\small},
  tail/.style={draw=ztgreen, rounded corners=1pt, thick, fill=ztgreen!8, minimum width=4.85cm, minimum height=0.72cm, align=center, font=\small},
  arrow/.style={->, thick}
]
  \node[title] (xt) at (-3.25,2.7) {x86\_64 trampoline};
  \node[hot] (x1) at (-3.25,1.75) {push probe id};
  \node[hot] (x2) at (-3.25,0.85) {call entry\_stub};
  \node[cell] (x3) at (-3.25,-0.05) {搬运后的原函数前导指令};
  \node[tail] (x4) at (-3.25,-0.95) {jmp func + orig\_len};

  \node[title] (at) at (3.25,2.7) {aarch64 trampoline};
  \node[hot] (a1) at (3.25,1.75) {mov x15, x30};
  \node[hot] (a2) at (3.25,0.85) {x17 = probe id，blr entry\_stub};
  \node[hot] (a3) at (3.25,-0.05) {mov x30, x15};
  \node[cell] (a4) at (3.25,-0.95) {搬运后的原函数前导指令};
  \node[tail] (a5) at (3.25,-1.85) {br func + orig\_len};

  \foreach \from/\to in {xt/x1,x1/x2,x2/x3,x3/x4,at/a1,a1/a2,a2/a3,a3/a4,a4/a5} {
    \draw[arrow] (\from) -- (\to);
  }
  \draw[dashed, thick, ztline] (0,-2.25) -- (0,3.1);
\end{tikzpicture}
}
\caption{两套 trampoline 内部组织对比}
\label{fig:trampoline-layout-compare}
\end{figure}

\begin{table}[H]
\centering
\caption{两套 stub 的现场保存重点}
\label{tab:stub-state-compare}
\begin{tabularx}{\textwidth}{P{3.1cm}Y Y}
\toprule
\textbf{现场类型} & \textbf{x86\_64} & \textbf{aarch64}\\
\midrule
整数参数 & 保存并恢复 SysV ABI 参数寄存器，payload 读取统一 \code{gp\_arg0..5}。 & 保存并恢复 \code{x0..x5}，payload 同样读取 \code{gp\_arg0..5}。\\
返回值 & exit stub 读取 \code{rax} 和 \code{xmm0} 对应槽位。 & exit stub 读取 \code{x0} 和 \code{q0} 对应槽位。\\
标志位 & 保存/恢复 \code{rflags}，避免条件判断受到 probe 影响。 & 保存/恢复 \code{nzcv}，保持条件码语义。\\
浮点状态 & 保存用于参数展示和返回值展示的 \code{xmm0..xmm7} 槽位。 & 保存用于 AAPCS64 参数展示和返回值展示的 \code{q0..q7} 槽位，并记录 \code{fpcr/fpsr}。\\
返回链状态 & TLS shadow stack 保存原返回地址，exit stub 弹出后 \code{ret}。 & TLS shadow stack 保存真实 LR，exit stub 弹出后写回 \code{x30} 并 \code{ret}。\\
\bottomrule
\end{tabularx}
\end{table}

\begin{table}[H]
\centering
\caption{前导指令重定位工作量对比}
\label{tab:relocation-compare}
\begin{tabularx}{\textwidth}{P{3.25cm}Y Y}
\toprule
\textbf{重定位对象} & \textbf{x86\_64 处理方式} & \textbf{aarch64 处理方式}\\
\midrule
普通指令 & 指令边界由 Capstone 给出，原样复制到 trampoline。 & 固定 4 字节指令，原样复制到 trampoline。\\
PC 相对取址 & RIP 相对内存访问需要按新旧 next IP 重新计算 \code{disp32}。 & \code{adr/adrp} 直接改写为等价的绝对地址装载序列。\\
直接调用 & \code{call rel32} 改写为绝对调用模板，避免 trampoline 与原目标距离过远。 & \code{bl imm26} 改写为绝对调用模板，目标地址由临时寄存器承载。\\
无条件跳转 & \code{jmp rel8/rel32} 改写为绝对跳转。 & \code{b imm26} 改写为绝对跳转。\\
条件跳转 & \code{jcc} 展开为条件反转和绝对跳转块。 & \code{b.cond/cbz/cbnz/tbz/tbnz} 改写为条件反转后进入绝对跳转块。\\
literal load & x86\_64 没有同类固定模式，主要落在 RIP 相对访存处理。 & \code{ldr literal} 的目标依赖 PC，当前安装阶段返回失败，避免生成语义不确定的 trampoline。\\
\bottomrule
\end{tabularx}
\end{table}

aarch64 后端使用 \code{PTRACE\_GETREGSET}/\code{PTRACE\_SETREGSET} 读写 \code{struct user\_pt\_regs}，远程 syscall 走 \code{x0..x5/x8}，远程函数调用使用 \code{x0/x1} 传参并从 \code{x0} 读取返回值。trampoline builder 已覆盖普通指令、\code{adr/adrp}、直接分支、条件分支、\code{cbz/cbnz} 和 \code{tbz/tbnz}。遇到尚未覆盖的 PC-relative 模式时，安装过程返回失败，tracer 保持原函数入口不变。

\begin{table}[H]
\centering
\caption{aarch64 后端关键约定}
\label{tab:aarch64-backend}
\begin{tabularx}{\textwidth}{P{3.6cm}Y}
\toprule
\textbf{项目} & \textbf{实现方式}\\
\midrule
入口 patch & 固定 16 字节模板：\code{ldr x16, \#8; br x16; .quad trampoline\_addr}。\\
远程 syscall & 临时写入 \code{svc \#0; brk \#0}，参数放入 \code{x0..x5}，系统调用号放入 \code{x8}，返回值从 \code{x0} 读取。\\
远程函数调用 & 临时 stub 使用 \code{ldr x16, \#16; blr x16; brk \#0}，当前用于 \code{x0/x1} 两个参数的目标进程内函数调用。\\
return probe & trampoline 先把原 \code{x30/lr} 放入 \code{x15}；entry stub 记录真实 LR 后把 \code{exit\_stub} 地址通过 \code{x15} 带回；trampoline 再写入 \code{x30}，函数返回时进入 exit stub。\\
参数与返回值 & 整数参数映射到 \code{x0..x5}，整数返回值映射到 \code{x0}；FP/SIMD 参数槽映射到 \code{q0..q7} 的低位快照，格式化层按统一 \code{fp\_args[]} 读取。\\
状态字段 & stub 记录并恢复 \code{nzcv}，同时保存 \code{fpcr/fpsr}，使整数条件码和浮点控制状态不会被 payload 调用路径改变。\\
\bottomrule
\end{tabularx}
\end{table}

\section{实验设计与结果}

\subsection{实验环境}

本节数据来自 x86\_64 开发机和 Google Cloud T2A aarch64 云服务器上的 benchmark 报告。
两组实验均通过 Makefile 的 benchmark 目标运行，脚本检测到 tracingfs 中的 uprobe 配置入口，因此包含 kernel uprobe 对照数据。
ARM 测试机为 Google Cloud \code{t2a-standard-4}，实例区域 \code{asia-southeast1-b}，CPU 为 4 核 ARM Neoverse-N1，系统为 Ubuntu 24.04.4 LTS，内核版本为 \code{6.17.0-1020-gcp}。

\begin{table}[H]
\centering
\caption{实验环境}
\label{tab:env}
\begin{tabularx}{\textwidth}{P{3.2cm}Y Y}
\toprule
\textbf{项目} & \textbf{x86\_64 开发机} & \textbf{aarch64 云服务器}\\
\midrule
机器规格 & 13th Gen Intel(R) Core(TM) i5-13400，16 logical CPUs & Google Cloud \code{t2a-standard-4}，4 vCPU，16 GiB RAM\\
CPU/架构 & Intel x86\_64 & ARM Neoverse-N1，aarch64\\
操作系统 & Linux 5.10.240-zeroVisor & Ubuntu 24.04.4 LTS，Linux 6.17.0-1020-gcp\\
实例区域 & 本地开发机 & Google Cloud \code{asia-southeast1-b}\\
编译器 & 系统默认 GCC & GCC 13.3.0\\
编译入口 & \code{make} & \code{make ARCH=aarch64}\\
性能入口 & \code{make benchmark} & \code{make benchmark}\\
LaTeX 编译入口 & \code{make paper} & \code{make paper}\\
\bottomrule
\end{tabularx}
\end{table}

\subsection{正确性实验}

正确性实验按赛题要求拆分。测试除了检查退出码，也记录关键状态变化，例如 entry/return 配平、条件过滤计数、热更新阶段范围、trampoline slot 复用和清理后的 \code{maps} 差分。

\begin{longtable}{@{}P{5.0cm}P{2.8cm}P{7.35cm}@{}}
\caption{自动化测试记录}\label{tab:test-coverage}\\
\toprule
\textbf{测试} & \textbf{状态} & \textbf{关键检查项}\\
\midrule
\endfirsthead
\toprule
\textbf{测试} & \textbf{状态} & \textbf{关键检查项}\\
\midrule
\endhead
\code{test\_libc\_trace} & \doneexp & 动态库函数、参数、返回值、字符串/buffer、format string 展开。\\
\code{test\_context\_integrity} & \doneexp & 通用寄存器、标志寄存器、浮点参数和浮点返回值上下文完整性。\\
\code{test\_probe\_lifecycle} & \doneexp & 16 probe 共存、条件过滤、call action、slot 分配与清理恢复。\\
\code{test\_thread\_safety} & \doneexp & 多线程并发命中、entry/return 配平、运行中 enable/disable。\\
\begin{tabular}[t]{@{}l@{}}\code{test\_thread\_}\\\code{group\_control}\end{tabular} & \doneexp & 新线程纳入 ptrace 控制、线程组 stop/continue 收敛。\\
\code{test\_probe\_hot\_update} & \doneexp & filter、call action、enable/disable 热更新。\\
\code{test\_poll\_exit\_race} & \doneexp & 短生命周期进程退出窗口，poll 正常收尾。\\
\code{test\_trace\_buffer\_reader} & \doneexp & ring buffer wrap、lost event、未提交 slot 处理。\\
\begin{tabular}[t]{@{}l@{}}\code{scripts/}\\\code{check\_arch\_config.py}\end{tabular} & \doneexp & x86\_64/aarch64 Makefile 后端选择检查。\\
\begin{tabular}[t]{@{}l@{}}\code{scripts/}\\\code{merge\_trace\_events.py}\\\code{--self-test}\end{tabular} & \doneexp & zeroTrace 与 perf/ftrace 风格日志排序合流。\\
\code{scan-build} & \doneexp & \code{scan-build --status-bugs make clean all} 已记录 No bugs found。\\
\bottomrule
\end{longtable}

测试源码按职责拆分。自动化 runner 放在 \code{cases}，被 trace 的目标程序放在 \code{targets}；benchmark、手动 demo 和共享工具也分别放在独立目录。基础生命周期、热更新、线程组控制和多线程运行时行为都有对应用例。职责拆开以后，失败日志更容易定位到具体机制。

测试记录覆盖了 F1--F7；A1 中的 x86\_64 后端已经跑完整套测试，aarch64 后端已有构建选择检查、trampoline builder 自测和 Google Cloud T2A 上的 benchmark 运行记录。

资源清理和退出窗口有单独压力记录。probe lifecycle 的 cleanup 子测试会在目标仍存活时记录 trampoline 地址、入口原始字节和 runtime mmap 区域；删除最后一个 probe 后，再检查进程 maps 和入口字节是否恢复。

目标退出窗口测试覆盖短生命周期进程退出、远程读取失败和 trace buffer 映射消失之间的竞态，已记录 1000 轮压力通过。trace buffer reader 测试会故意延迟首次 poll，让事件数量超过 ring buffer 容量，要求日志显式出现 \code{ztrace:lost}，把被覆盖事件的范围写出来。

\subsection{性能实验方法}

性能实验包含四项：baseline、zeroTrace、uprobe 和生命周期延迟。baseline 直接循环调用 \code{bench\_getpid()}；zeroTrace 在同一目标函数上安装 probe 后执行相同循环；uprobe 在权限满足时使用 bpftrace 对同一函数采样；生命周期延迟对同一个 probe 执行 1000 轮 install/uninstall。

benchmark 目标函数是一个 noinline 的 getpid 包装函数，内部直接发起 getpid syscall。这样可以避开 libc/vDSO 的实现差异，让 baseline、kernel uprobe 和 zeroTrace 三组实验尽量测量同一个用户态函数入口。

默认参数为 1,000,000 次调用、5 轮重复，脚本输出 mean/min/max/stdev。kernel uprobe 对照依赖 \code{bpftrace} 与 tracingfs 权限；权限满足时纳入报告，权限不足时脚本会写出说明并跳过该项。

核心指标定义为：
\begin{equation}
\begin{aligned}
  overhead\_per\_call &=
  \frac{ztrace\_total\_ns - baseline\_total\_ns}{iterations},\\
  install\_latency &=
  \frac{\sum_{i=1}^{N} install\_ns_i}{N}, \quad
  uninstall\_latency =
  \frac{\sum_{i=1}^{N} uninstall\_ns_i}{N}.
\end{aligned}
\end{equation}

\subsection{性能结果}

\begin{longtable}{@{}P{4.9cm}P{5.0cm}P{5.0cm}@{}}
\caption{x86\_64 与 aarch64 benchmark 结果}
\label{tab:benchmark}\\
\toprule
\textbf{指标} & \textbf{x86\_64} & \textbf{aarch64 / T2A}\\
\midrule
\endfirsthead
\toprule
\textbf{指标} & \textbf{x86\_64} & \textbf{aarch64 / T2A}\\
\midrule
\endhead
iterations & 1,000,000 & 1,000,000\\
repeats & 5 & 5\\
baseline total ns mean & 61,239,703 ns & 173,445,186 ns\\
baseline per call mean & 61.24 ns & 173.45 ns\\
baseline per call min/max & 60.59 ns / 62.09 ns & 173.11 ns / 173.67 ns\\
baseline per call stdev & 0.66 ns & 0.28 ns\\
uprobe total ns mean & 1,994,438,922 ns & 653,140,257 ns\\
uprobe per call mean & 1,994.44 ns & 653.14 ns\\
uprobe per call min/max & 1,905.74 ns / 2,052.76 ns & 650.77 ns / 657.08 ns\\
uprobe per call stdev & 58.88 ns & 2.53 ns\\
uprobe overhead/call mean & 1,933.20 ns & 479.70 ns\\
uprobe overhead/call min/max & 1,844.99 ns / 1,990.67 ns & 477.11 ns / 483.97 ns\\
uprobe overhead/call stdev & 58.33 ns & 2.67 ns\\
ztrace total ns mean & 235,684,856 ns & 410,467,298 ns\\
ztrace per call mean & 235.68 ns & 410.47 ns\\
ztrace per call min/max & 233.38 ns / 239.31 ns & 391.45 ns / 420.28 ns\\
ztrace per call stdev & 2.55 ns & 11.18 ns\\
ztrace overhead/call mean & 174.45 ns & 237.02 ns\\
ztrace overhead/call min/max & 172.63 ns / 177.53 ns & 217.84 ns / 246.61 ns\\
ztrace overhead/call stdev & 2.28 ns & 11.25 ns\\
ztrace vs uprobe & 11.08x lower overhead & 2.02x lower overhead\\
install latency avg & 370,539 ns，约 0.371 ms & 599,175 ns，约 0.599 ms\\
uninstall latency avg & 83,892 ns，约 0.084 ms & 215,881 ns，约 0.216 ms\\
\bottomrule
\end{longtable}

\begin{figure}[H]
\centering
\begin{tikzpicture}[
  x=1cm,
  y=1cm,
  axis/.style={thick},
  bar/.style={draw=black, thick, rounded corners=1pt}
]
  \path[use as bounding box] (-1.1,-0.75) rectangle (12.2,5.25);

  \node[font=\small] at (2.3,5.0) {x86\_64 函数调用耗时};
  \draw[axis] (0,0) -- (4.6,0);
  \draw[axis] (0,0) -- (0,4.5);
  \foreach \value/\pos in {0/0,500/1.125,1000/2.25,1500/3.375,2000/4.5} {
    \draw (-0.05,\pos) -- (0.05,\pos) node[left] {\scriptsize \value};
  }
  \node[rotate=90] at (-0.85,2.25) {\small ns/call};

  \draw[bar, fill=ztgreen!35] (0.75,0) rectangle (1.35,0.14);
  \node[above] at (1.05,0.14) {\scriptsize 61.24};
  \node[below] at (1.05,-0.18) {\scriptsize baseline};

  \draw[bar, fill=ztblue!35] (2.0,0) rectangle (2.6,0.53);
  \node[above] at (2.3,0.53) {\scriptsize 235.68};
  \node[below] at (2.3,-0.18) {\scriptsize ztrace};

  \draw[bar, fill=ztorange!45] (3.25,0) rectangle (3.85,4.49);
  \node[above] at (3.55,4.49) {\scriptsize 1994.44};
  \node[below] at (3.55,-0.18) {\scriptsize uprobe};

  \node[font=\small] at (9.0,5.0) {x86\_64 probe 生命周期延迟};
  \draw[axis] (6.8,0) -- (11.2,0);
  \draw[axis] (6.8,0) -- (6.8,4.5);
  \foreach \label/\pos in {0/0,0.1/1.125,0.2/2.25,0.3/3.375,0.4/4.5} {
    \draw (6.75,\pos) -- (6.85,\pos) node[left] {\scriptsize \label};
  }
  \node[rotate=90] at (6.0,2.25) {\small ms};

  \draw[bar, fill=ztblue!35] (8.0,0) rectangle (8.6,4.17);
  \node[above] at (8.3,4.17) {\scriptsize 0.371};
  \node[below] at (8.3,-0.18) {\scriptsize install};

  \draw[bar, fill=ztgreen!35] (9.75,0) rectangle (10.35,0.95);
  \node[above] at (10.05,0.95) {\scriptsize 0.084};
  \node[below] at (10.05,-0.18) {\scriptsize uninstall};
\end{tikzpicture}
\caption{x86\_64 benchmark 柱状图：函数调用开销与 probe 生命周期延迟}
\label{fig:perf-bars-x86}
\end{figure}

\begin{figure}[H]
\centering
\begin{tikzpicture}[
  x=1cm,
  y=1cm,
  axis/.style={thick},
  bar/.style={draw=black, thick, rounded corners=1pt}
]
  \path[use as bounding box] (-1.1,-0.75) rectangle (12.2,5.25);

  \node[font=\small] at (2.3,5.0) {aarch64 函数调用耗时};
  \draw[axis] (0,0) -- (4.6,0);
  \draw[axis] (0,0) -- (0,4.5);
  \foreach \value/\pos in {0/0,200/1.286,400/2.571,600/3.857} {
    \draw (-0.05,\pos) -- (0.05,\pos) node[left] {\scriptsize \value};
  }
  \node[rotate=90] at (-0.85,2.25) {\small ns/call};

  \draw[bar, fill=ztgreen!35] (0.75,0) rectangle (1.35,1.12);
  \node[above] at (1.05,1.12) {\scriptsize 173.45};
  \node[below] at (1.05,-0.18) {\scriptsize baseline};

  \draw[bar, fill=ztblue!35] (2.0,0) rectangle (2.6,2.64);
  \node[above] at (2.3,2.64) {\scriptsize 410.47};
  \node[below] at (2.3,-0.18) {\scriptsize ztrace};

  \draw[bar, fill=ztorange!45] (3.25,0) rectangle (3.85,4.20);
  \node[above] at (3.55,4.20) {\scriptsize 653.14};
  \node[below] at (3.55,-0.18) {\scriptsize uprobe};

  \node[font=\small] at (9.0,5.0) {aarch64 probe 生命周期延迟};
  \draw[axis] (6.8,0) -- (11.2,0);
  \draw[axis] (6.8,0) -- (6.8,4.5);
  \foreach \label/\pos in {0/0,0.2/1.286,0.4/2.571,0.6/3.857} {
    \draw (6.75,\pos) -- (6.85,\pos) node[left] {\scriptsize \label};
  }
  \node[rotate=90] at (6.0,2.25) {\small ms};

  \draw[bar, fill=ztblue!35] (8.0,0) rectangle (8.6,3.85);
  \node[above] at (8.3,3.85) {\scriptsize 0.599};
  \node[below] at (8.3,-0.18) {\scriptsize install};

  \draw[bar, fill=ztgreen!35] (9.75,0) rectangle (10.35,1.39);
  \node[above] at (10.05,1.39) {\scriptsize 0.216};
  \node[below] at (10.05,-0.18) {\scriptsize uninstall};
\end{tikzpicture}
\caption{Google Cloud T2A aarch64 benchmark 柱状图：函数调用开销与 probe 生命周期延迟}
\label{fig:perf-bars-aarch64}
\end{figure}

在 x86\_64 结果中，zeroTrace 额外开销约 174.45 ns/call，低于赛题 1000 ns 指标，也低于项目内部设定的 200 ns 优化目标。Google Cloud T2A aarch64 结果中，zeroTrace 额外开销约 237.02 ns/call，kernel uprobe 额外开销均值约 479.70 ns/call，zeroTrace 为 2.02x lower overhead。两套环境中的 install/uninstall 延迟均低于 1 ms，距离赛题 10 ms 指标还有较大余量。

\subsection{多线程与热更新压力结果}

多线程测试记录 entry/return 配平、线程纳管和运行中切换情况。已记录的 10 轮压力摘要如下：

\begin{lstlisting}
thread_safety: 10/10 passed
  tids=12 tracked=13 toggles=12 in every round
  thread_stable strict pairs: 2338..2400 per round
  thread_pair strict pairs  : 471..480 per round
  thread_mix high-frequency : >=21039 entry/return events per side

thread_group_control: 10/10 passed
  rounds=80 resume_rounds=80 in every round
  task_max/tracked_max range: 31..32, matched in every round
\end{lstlisting}

热更新测试检查 filter A/B/final、disable 阶段零泄漏、call action A/B 切换和 live call action 更新。已记录 100 轮压力结果：

\begin{lstlisting}
test_probe_hot_update: 100/100 passed
last live call action round: hot_call_a=170 hot_call_b=45
\end{lstlisting}

这些压力结果对应并发命中、运行中状态切换、返回链配平和 call action 半写入窗口，比单次手动演示更接近真实使用场景。A4 的测试还覆盖了带参 call action：测试把 entry 参数和常量一起传给目标进程内的 callee，并检查不同参数下的返回值。slot 碰撞子测试会保留一个旧 probe，同时反复 trace/untrace 推进新 probe id，覆盖两个 probe 映射到同一 action 槽位的场景，要求两者的 call action 同时保持正确。

\section{使用与复现}

\subsection{依赖与运行前准备}

项目依赖 \code{libcapstone}、\code{libdl} 和 \code{libreadline}。在 Debian/Ubuntu 环境中可以使用以下命令安装构建依赖：

\begin{lstlisting}
sudo apt install build-essential libcapstone-dev libreadline-dev
\end{lstlisting}

zeroTrace 依赖 ptrace 控制目标进程。现代 Linux 发行版通常通过 Yama 限制 ptrace 附加其他进程，实验前需要确认 \code{/proc/sys/kernel/yama/ptrace\_scope}。若系统策略允许，可临时设置为 0：

\begin{lstlisting}
cat /proc/sys/kernel/yama/ptrace_scope
echo 0 | sudo tee /proc/sys/kernel/yama/ptrace_scope
\end{lstlisting}

\subsection{构建与测试}

项目提供统一 Makefile。常规构建、测试和 benchmark 命令如下：

\begin{lstlisting}
make
make test
make benchmark
\end{lstlisting}

Makefile 会根据构建机器选择架构后端，也可以显式指定：

\begin{lstlisting}
make ARCH=x86_64
make ARCH=aarch64
make ARCH=aarch64 CC=aarch64-linux-gnu-gcc
\end{lstlisting}

LaTeX 文档可以本地编译：

\begin{lstlisting}
make paper
make clean-paper
\end{lstlisting}

\code{paper} 目标使用 \code{latexmk -xelatex}，文档使用 \code{ctexart} 的 Fandol 字体集，Linux 环境可以直接编译。

\subsection{CLI 示例}

zeroTrace 的交互方式接近调试器。用户先 attach 目标进程，再安装、更新或卸载 probe。

\begin{lstlisting}
./bin/ztrace
attach <pid>
trace add_loop
trace write if arg0 == 1 && arg2 > 0
update add_loop if arg0 >= 100
update add_loop call call_marker_args arg0 0x10 arg0
disable add_loop
enable add_loop
untrace add_loop
detach
\end{lstlisting}

日志默认写入工作目录下的 \code{ztrace.<pid>.log}。如果 \code{conf/zttrace.conf} 中存在函数签名，日志会优先输出接近函数调用形式的内容，例如：

\begin{lstlisting}
test_loop-532386/532386 [010] 58915.513461172:
  ztrace:entry: add_loop(24, 25)
test_loop-532386/532386 [010] 58915.513468205:
  ztrace:return: add_loop -> 49
test_loop-532386/532386 [010] 58915.513461172:
  ztrace:entry: fp_add_loop(24.25, 1.5)
test_loop-532386/532386 [010] 58915.513468205:
  ztrace:return: fp_add_loop -> 25.75
\end{lstlisting}

\section{分工说明}

\begin{table}[H]
\centering
\caption{队员分工}
\begin{tabularx}{\textwidth}{P{2.4cm}Y}
\toprule
\textbf{成员} & \textbf{主要工作}\\
\midrule
舒钰尧 & 主要负责 trace runner 与探针生命周期管理，完成 attach、trace、untrace、enable、disable、update 等上层流程串联；参与 return probe 调用配对设计，整理 entry event、return event 与 call event 的状态关系；负责部分自动化测试用例编写与回归验证。\\
曹佳一 & 主要负责 injector 与目标进程控制相关工作，包括 ptrace 线程组控制、远程内存读写、远程 mmap/munmap、payload 注入、入口 patch 与原始指令恢复；重点处理多线程目标进程中的线程枚举、停止、继续和 patch 安全窗口。\\
倪晨曦 & 主要负责 trampoline、stub 与 ISA 后端实现，完成 x86\_64/aarch64 入口跳转模板、前导指令保存与跳回逻辑；参与 payload 侧现场保存、参数/返回值捕获、条件探针和 call action 联调。\\
\bottomrule
\end{tabularx}
\end{table}

\section{AI 辅助工具使用说明}

开发和文档整理过程中使用过 AI 辅助工具，主要用于讨论改进方案、辅助编写重复性或规则明确的代码、整理测试清单和调整文档表述。项目架构、trampoline/stub/payload 机制、线程组控制策略、性能优化方向和最终取舍由团队完成；报告中的代码、测试和性能结论以仓库内容和本地运行输出为准。

\begin{table}[H]
\centering
\caption{AI 辅助工具使用说明}
\label{tab:ai-usage}
\begin{tabularx}{\textwidth}{P{3.0cm}Y}
\toprule
\textbf{字段} & \textbf{说明}\\
\midrule
工具名称 & OpenAI ChatGPT / Codex 类代码助手\\
模型名称 & OpenAI GPT-5\\
使用范围 & 改进方案讨论、重复性代码辅助、边界测试建议、benchmark 现象分析和文档措辞整理\\
人工完成部分 & 总体架构、关键机制设计、核心实现取舍、实验设计和最终结果解释\\
人工校验 & diff 审阅、构建、定向测试、全量测试、benchmark 和文档事实核对\\
最终责任 & 团队负责方案确认、代码审阅、测试验证和最终提交\\
\bottomrule
\end{tabularx}
\end{table}

\section{结论}

zeroTrace 已经打通从安装 probe 到输出日志的完整流程：tracer 负责安装和管理，目标进程内部完成命中后的采集逻辑，事件通过共享 trace buffer 异步传出。项目实现了基础功能 F1--F7，并补充了条件探针、热更新、探针内调用、perf/ftrace 风格日志和多架构后端。x86\_64 benchmark 显示，单次 probe 额外开销保持在 200 ns 以内；Google Cloud T2A aarch64 benchmark 中，zeroTrace 额外开销均值为 237.02 ns/call，相对 kernel uprobe 为 2.02x lower overhead。两套环境中的安装/卸载延迟均远低于 10 ms。

\begin{thebibliography}{9}
\bibitem{linux-ptrace}
Linux manual page: ptrace(2).
\bibitem{linux-process-vm-readv}
Linux manual page: process\_vm\_readv(2).
\bibitem{perf-script}
Linux perf documentation: perf-script.
\bibitem{ltrace-conf}
ltrace configuration manual: ltrace.conf(5).
\bibitem{nemu}
NJU ProjectN, NEMU expression evaluator.
\end{thebibliography}

\end{document}
